\section{Equações Lineares Parabólicas de 2ª ordem}


\subsection*{Funções em Espaços de Banach}
\begin{frame}{Funções em Espaços de Banach}
Seja \(T>0\) um número real e {\color{blue}\((X,\|\cdot\|)\)}	um espaço de Banach real. Para uma função {\color{blue}\(u:[0,T]\longrightarrow X\)} temos os seguintes conceitos, cujo enunciado preciso podem ser encontrados em \cite[\S E.5]{evans},
\begin{itemize}
\item {\color{blue}\(u\)} é dita ser {\color{blue}fortemente mensurável} quando é o limite de uma sequências de funções simples em \(X\).
\item  {\color{blue}\(u\)} é dita ser {\color{blue} somável ou integrável} quando é fortemente mensurável e a função real \(t\longmapsto {\color{blue}\|u(t)\|}\) é integrável. Neste caso, a sua integral é um elemento de {\color{blue}\(X\)} e será denotado por:
\[\int_0^T {\color{blue}u(t)}\,dt.\]
\end{itemize}
\end{frame}
%%%%%%%%%%%%%%%%%%%%%%%%%%%%%%%%%%%%%%%%%%%%%%%%%%

\begin{frame}
\begin{teo}
Se \(u:[0,T]\longrightarrow X\) é integrável, então
\[
\left\|\int_0^T u(t)\,dt\right\|\leq \int_0^T \|u(t)\|\,dt
\]
e, para toda \(A\in X'\), 
\[
\left\langle A, \int_0^T u(t)\,dt\right\rangle_{X'X}=\int_0^T \left\langle A, u(t)\right\rangle_{X',X}\, dt
\]
\end{teo}
\end{frame}

%%%%%%%%%%%%%%%%%%%%%%%%%%%%%%%%%%%%%%%%%%%%%%%%%%

\begin{frame}
\begin{defin}
O espaço \(L^p(0,T;X)\) consiste das funções \(u:[0,T]\longrightarrow X\) integráveis tais que \(\n{u(\cdot)}\in L^p(0,T)\). Em \(L^p(0,T;X)\), definimos a norma:
\[
\n{u}_{L^p(0,T;X)}=\left(\int_0^T \n{u(t)}^p\right)^{1/p},\ 1\leq p<\infty, 
\]
e
\[
\n{u}_{L^\infty(0,T;X)}=\esup_{0\leq t\leq T}\n{u(t)}.
\]
\end{defin}

\begin{exer}
Mostre que \(L^p(0,T;X)\) é um espaço de Banach.
\end{exer}

\subsection*{Exercício \theexer}

\end{frame}